\chapter{Background: Genetic Programming}
\label{ch:background}

\section{Evolutionary Algorithms}
Genetic algorithms evolve a population of candidate solutions with biologically
motivated operators. They need no gradient information, use probabilistic
rather than deterministic search, and work well on non-linear, non-convex and
multimodal problems. A generic run has these components: a genetic
representation of a solution (the \emph{genome}), a way to create new solutions,
a \emph{fitness function}, a \emph{selection} function, \emph{crossover} and
\emph{mutation}. The set of solutions at a given step is a \emph{generation}.

As an example, in the knapsack problem a genome is a 0/1 array saying which items
are packed, the fitness is the total value if the weight limit is respected
(and 0 otherwise), single-point crossover swaps the tails of two parents, and
\emph{elitism} copies the best solutions unchanged into the next generation so
they are not destroyed by crossover.

\section{Genetic Programming}
Genetic Programming \citep{koza1994gp} applies the same loop to \emph{programs},
usually represented as trees: internal nodes are functions and leaves are
terminals. Fitness is measured by running the program on the problem.

\subsection{Initialisation: ramped half-and-half}
Half of the trees are built with the \emph{full} method (every leaf at the
maximum depth), the other half with the \emph{grow} method (random shapes within
the maximum depth), with the depth ramped across the population to obtain
diverse sizes and shapes.

\subsection{Selection, crossover, mutation}
\emph{Tournament selection} draws $k$ individuals and keeps the fittest.
\emph{Subtree crossover} swaps random subtrees between two parents.
\emph{Mutation} perturbs a node or replaces a subtree with a new random one.
\emph{Hoist mutation} replaces a tree by one of its own subtrees and is used to
fight \emph{bloat}, the uncontrolled growth of trees.

\section{Genetic Programming Hyper-Heuristics}
In a hyper-heuristic, GP does not solve the problem directly; it evolves the
\emph{heuristic} (here a priority function) that then solves instances of the
problem. The evolved rule is evaluated on training instances and should
generalise to unseen ones.
